% ============================================================
% FUENTES: draft p.1--2 (transcripción ## Page 1--2).
% ============================================================
\section{Sobre la evaluación y el temario}

\subsection{Reglas del juego y esquema de evaluación}

Bienvenidos al curso de Matemáticas Biológicas. Antes de adentrarnos formalmente en las ecuaciones diferenciales y en la dinámica no lineal, dejemos claras las reglas de evaluación y la dinámica de trabajo:

\begin{itemize}
    \item \textbf{Mini Proyectos:} Tareas-proyecto extensas de desarrollo teórico y numérico. Usualmente se asignan con un plazo de 15 días para su entrega y resolución cuidadosa.
    \item \textbf{Mini Labs:} Prácticas computacionales opcionales, pero sumamente útiles para perderle el miedo a la programación y a la simulación numérica de sistemas dinámicos.
    \item \textbf{Proyecto Final:} Se desarrollará durante el último cuarto del semestre (aproximadamente en el último mes) y consiste en la modelación y análisis cualitativo completo de un fenómeno biológico de interés.
\end{itemize}

\paragraph{Política sobre citas, copiado y plagio.}
Sé bien cómo funcionan estas dinámicas y sé que van a consultar y discutir material entre ustedes. Los problemas del curso han sido diseñados específicamente para estas notas, por lo que la gran mayoría no tiene soluciones directas flotando en internet. Si colaboran o toman una idea de un compañero o de un libro, no me molesta en absoluto, \emph{siempre y cuando incluyan la referencia explícita y den el crédito correspondiente}. Lo que no tolero es el plagio deshonesto (presentar ideas ajenas como propias sin citar).

\aside{Una advertencia fundamental: lo que bajo ninguna circunstancia se aceptará es que entreguen fotos o capturas de pantalla del trabajo de otra persona. Al menos, usar las propias manos para transcribir, escribir en \LaTeX{} o teclear el código garantiza que procesen e incorporen los conceptos.}

\subsection{Filosofía del curso y precedentes necesarios}

Los cursos de modelación biológica y sistemas dinámicos son demandantes porque requieren conectar herramientas de diversas ramas de las matemáticas:
\begin{enumerate}
    \item Ecuaciones diferenciales ordinarias (EDOs).
    \item Álgebra lineal (valores propios, transformaciones lineales).
    \item Cálculo diferencial e integral (esencialmente saber graficar curvas, raíces e intersecciones).
    \item Pensamiento geométrico y cualitativo.
    \item Cierto sentido común o ``conocimiento biológico''\aside{Sea lo que sea que eso signifique en la práctica transdisciplinaria.} para interpretar variables y parámetros con significado biológico concreto.
\end{enumerate}

Si han cursado materias de otras disciplinas o tienen interés en la transdisciplina, este marco les resultará natural y divertido. La modelación matemática consiste en aprender a jugar con las herramientas analíticas para develar la estructura oculta de los fenómenos naturales.

\paragraph{Un enfoque geométrico, no axiomático-formalista.}
Este curso no seguirá la estructura clásica de \emph{Definición $\to$ Teorema $\to$ Demostración $\to$ Corolario}. En lugar de perdernos en tecnicismos analíticos que a menudo ocultan la física del problema, nos apoyaremos en intuición cualitativa, argumentos geométricos, retratos de fase y muchos diagramas.

\aside{No se dejen engañar por la aparente sencillez visual de los primeros temas: si descuidan los detalles y las sutilezas geométricas, el curso puede deparar sorpresas mayúsculas al avanzar hacia bifurcaciones y caos.}

Para construir estos cimientos, comenzaremos repasando la evolución histórica que llevó a la creación de la dinámica no lineal moderna.

